\section{总结与延伸}

\subsection{核心论点}

\begin{enumerate}
    \item \textbf{从 Scaling 时代回到 Research 时代}：Pre-training 的数据墙和 RL 的泛化局限意味着纯粹的规模扩张不再够用。我们需要根本性的新想法——但现在有了大型计算机来验证它们
    \item \textbf{泛化是根本挑战}：当前模型在 eval 上表现出色但真实能力不足，根源在于泛化能力远不如人类。人类在进化无法提供先验的领域（如编程）仍展现出卓越的学习效率
    \item \textbf{超级智能 = 类人持续学习 Agent}：不是"无所不知的 oracle"，而是"学什么都极快的 mind"。部署即学习，多实例可融合知识
    \item \textbf{安全需要行动，不仅是讨论}：增量部署、展示 AI 能力、竞争者之间合作、限制最强系统的力量——都是务实的安全策略
    \item \textbf{情感作为 Value Function}：进化编码高层级社会欲望的方式仍是未解之谜，理解它可能为 AI 对齐提供关键线索
    \item \textbf{Ideas 比 Execution 更稀缺}：在所有人做同样事情的 scaling 时代，差异化的想法成为最稀缺的资源
\end{enumerate}

\subsection{把访谈压缩成一张研究决策表}

\begin{table}[H]
\centering
\begin{tabular}{p{3.1cm}p{3.8cm}p{4.6cm}}
\toprule
\textbf{议题} & \textbf{Ilya 的判断} & \textbf{对研究者的实际含义} \\
\midrule
Pre-training & 仍重要，但已接近数据墙 & 继续扩规模不再足够，必须寻找新范式 \\
RL & 代表新的学习机会，但泛化仍弱 & 需要更好的任务设计、验证方式与训练信号 \\
Superintelligence & 更像持续学习 agent，而非静态 oracle & 系统应强调在线学习、记忆和知识整合 \\
Safety & 需要渐进部署与组织级约束 & 研究不能只谈原则，还要设计治理机制 \\
Research taste & 好想法比执行更稀缺 & 研究者要培养美感、简洁性与 top-down belief \\
\bottomrule
\end{tabular}
\caption{如果把整场访谈用于研究决策，它更像一张路线选择表而不是单点观点合集}
\end{table}

\begin{importantbox}{给研究者最现实的提醒}
这场访谈最不浪漫也最有用的地方在于：Ilya 并没有说 ``scaling 结束了''，而是说 ``仅靠 scaling 已经不够''。对研究团队而言，这意味着未来最稀缺的资源既不是 GPU，也不是热点词，而是能把新想法变成可验证系统的判断力。
\end{importantbox}

\subsection{开放问题}

如果把这场访谈放回 2026 年的研究语境，还会留下几道没有被真正回答的问题：
\begin{itemize}
    \item 如果持续学习 agent 真是通往 superintelligence 的主线，那么什么样的 online data 和 feedback 才足以支撑它？
    \item RL 的泛化问题到底是 reward 设计问题、环境构造问题，还是模型结构本身的问题？
    \item 安全中的 ``限制最强系统权力'' 如何落成工程与治理机制，而不是停留在原则口号？
    \item 当 ideas 比 execution 更稀缺时，研究组织应该如何筛选值得长期下注的方向？
\end{itemize}

这些问题没有标准答案，但恰恰构成了 Ilya 所说的 ``Age of Research'' 的真正含义：不是把现有范式再做大一点，而是重新决定哪些问题值得被当作主问题。

\subsection{拓展阅读}

\begin{itemize}
    \item Kaplan et al., "Scaling Laws for Neural Language Models" (2020) --- Pre-training Scaling Laws
    \item SSI (Safe Superintelligence Inc.) 官网
    \item Dwarkesh Patel 关于 RL scaling sigmoid 曲线的博客文章
    \item Antonio Damasio, \textit{Descartes' Error} --- 情感在决策中的核心作用
    \item AlexNet, Transformer, ResNet 原始论文 --- 关键突破的计算规模对照
    \item OpenAI-Anthropic 安全合作公告 --- Ilya 预测的实例
\end{itemize}

%% --- Fin del contenido --- %%

\end{document}
